\chapter{R2 with both the native source ratio and relative current unbounded}
\label{component:pef021eabacde}
\noindent\textbf{Component }\path{pef021eabacde}\par
{\small\noindent Source: \path{NEW/root/r2_unbounded_current_rederived.tex}\par}
{\small\noindent SHA256: \path{ef021eabacde61ad9b1ed3c993a53db009b9974528f85eb04cfab31404534ea5}\par}
{\small\noindent Inventory use: UL1-UL18\par}
\medskip
% Fresh rederivation. This is not a recovered historical R21 component.
\section{R2 with both the native source ratio and relative current unbounded}
\label{pef021eabacde:sec:r2-unbounded-current-new}

\paragraph{Precise assertion.}
Retain the original RER source, all its strict gaps, and the notation
TU2--TU16 of the preceding bounded-relative-current proof.
Assume the native $a,t$ are positive unbounded and
\[
             \lambda=\frac{n}{tB}=\frac{n}{qz}
                   \text{ is positive unbounded}.                 \tag{UL1}
\]
Then the whole original weak-C2 parent is tensor SOS for every
independent projective test and every state. The argument below
complements TU1; it never assigns a sign to the formal C1 coefficient
$\beta^0$, and never replaces the original physical source by that
formal source.

\subsection{Source bounds without a thinness assumption}

The exact RER physicality equation already used in TU4 is
\[
 Q_{\rm RER}=\tau\{A+2r(z\Delta)\lambda\},\qquad
 \tau=\frac q{zw},\qquad
 A=(1+z)(1+Sh)(1+Sq).
\]
Here $A,Q_{\rm RER},r,z\Delta$ are positive units with residue one.
Since $\lambda$ is unbounded, this proves, with no cancellation of
mixed factors,
\[
 \tau\in\mathrm I,\qquad \lambda\tau\sim\frac12,\qquad
 \lambda q\sim\frac{zw}{2},\qquad
 \frac n{z^2w}\sim\frac12.                                     \tag{UL2}
\]
All four divisions are pure source. Recall $0<w<R\sim1$.
The exact TU9 formulas therefore imply
\[
\begin{gathered}
 (1-h)\delta c_b,\ h\delta a_b,\ \delta D_b,\
            \delta h_2,\delta h_b=\lambda q\,O(1),\\
 \delta\alpha_b/q^2=\lambda q\,O(S+zw),\qquad
 h_2\asymp zw,\quad h_b=zw/2,\\
 H_0=4h_2/R+2z\asymp z,\qquad H_0\geq2z,\\
 \frac g{\sqrt{H_0}}=O(1),\quad
 \frac c{\sqrt{H_0}}=O(1),\quad
 \frac{\lambda q}{H_0}=O(1),\quad \frac n{z^2}=O(1).
\end{gathered}                                                   \tag{UL3}
\]
In particular $\lambda q$ is infinitesimal. The displayed coefficient
arrays have denominators only among the pure-source units
$P,J,Q,R$ after the original $h,1-h$ openings cancel.
The actual and formal $\sigma=c_b+a_b$ differ by a relative
infinitesimal. Their square-root ratio is a bounded unit.
No lower bound on $w,h,1-h$ or the strict gap $M$ is assumed.

Use the complete generic SC3 identity with the actual coefficients.
Retain its positive $m$ packet and use $\widehat N=\sqrt{H_0}N$.
The formal comparison is used only for finite scalar and normal
identities. Its $q$-coquotients and the QI coefficient arrays are
bounded here: their derivation requires $S,z,q$ infinitesimal and
$w$ bounded, but not $\beta^0\geq0$. The sole QI bound that is no
longer available, $w/(c/g)=O(1)$, is replaced explicitly below.

\subsection{The entire extra current supplies a new angular pivot}

Use the exact positive-chart rows TU15--TU16, denoted here by
$A_*,B_*,C_*,D_*,U_*$ and $\Gamma$. Define
\[
 J_3=\sqrt\lambda(U_*-B_*),\qquad
 E_*=h_t-v,\qquad F_*=vh_t-gD_*.
                                                               \tag{UL4}
\]
The first line of the exact whole-current identity TU17 is
\[
 4D_uD_v\left(J_3-d_N\widehat N\right)^2,\qquad
 d_N=\frac{\sqrt\lambda\,cz}{\sqrt{H_0}}
       =\sqrt{\frac{nz}{H_0}}\in\mathrm I.                       \tag{UL5}
\]
Thus it supplies a strict real pivot on $J_3$, with only
infinitesimal changes to the hard normal Gram. Its scalar
remainder is still exactly
$8\lambda D_uD_v(1-e)\Gamma^2
-4\lambda D_uD_v(1-e^2)B_*^2$.
Its complete first-normal remainder, retained below, is
$-8\lambda D_uD_v(1-e)czB_*N$.

Write $\widetilde c,\widetilde a,\widetilde d,\widetilde\alpha,
\widetilde m,\widetilde b$ for the six differences
$\delta c_b,\delta a_b,\delta D_b,\delta\alpha_b,
\delta h_2,\delta h_b$ divided by $\lambda q$.
They are defined explicitly by the six right sides in TU9, using
$n=\lambda qz$; their apparent opening poles cancel in every
combination below. Put $U=u^2,V=v^2$, and retain
$\rho_0=Kh/R$, $a_t=1-\rho_0+(w/R)V$.
Define the exact rational arrays
\[
\begin{split}
\mathsf a={}&4D_vR^2\{\widetilde c(1-\rho_0)^2
                      +\widetilde a\rho_0^2/D_u\}
 +4\widetilde\alpha zS^2VR^2\rho_0^2/D_u
 +4\widetilde d S^2R^2a_t^2,\\
\mathsf b={}&-8\widetilde\alpha zSR\rho_0
                         -8\widetilde d SRa_tD_u,\\
\mathsf c={}&4\widetilde dD_u^2,\\
\mathsf l={}&4R\{\widetilde m(1-w/R)
                              +\widetilde b(w/R)D_v\}.
\end{split}                                                     \tag{UL6}
\]
For example $\delta a/q=\lambda\mathsf a$ in QI2,
and $\delta(R^2P_L/q)=\lambda\mathsf l$.
Direct specialization of TU9 gives, keeping $S,z,w,h$ exact,
\[
\begin{split}
 (\mathsf a-\mathsf c)|_{q=U=V=0}&=0,\\
 (\mathsf b+2\mathsf c+8D_uD_v(1-e))|_{q=U=V=0}&=0,\\
 (\mathsf l-4D_uD_v(1-e^2))|_{q=U=V=0}&=0.
\end{split}                                                     \tag{UL7}
\]
These are stronger than residue equalities. Clear the unit
denominators and divide successively by the monic variables
$q,U,V$. Each of the three expressions is consequently
$qF_q+UF_U+VF_V$, with finite bounded rational coefficient arrays.
The division does not invert $q,U,V$ in the tensor ring.

The combined full scalar correction to QI6 therefore consists of
bounded-array multiples of $\lambda q,\lambda U,\lambda V$
on $A_*^2,\Gamma^2,B_*^2$, an infinitesimal coefficient on
$E_*^2$, and the infinitesimal $\alpha$ correction on $D_*^2$.
The following exact identities allocate every apparently unbounded
testing term:
\[
\begin{aligned}
\lambda U A_*^2&=\lambda q(uE_*+bC_*)^2,\\
\lambda V A_*^2&=\lambda q(F_*+gD_*)^2,\\
\lambda U\Gamma^2&=\lambda q\,bC_*(uE_*+bC_*),\\
\lambda V\Gamma^2&=\lambda q(F_*+gD_*)
                                    (F_*+gD_*-vE_*),\\
\sqrt\lambda\,uB_*&=\sqrt{nS/z^2}\,(F_*+gD_*),\\
\sqrt\lambda\,vB_*&=\sqrt n\,C_*-vJ_3 .
\end{aligned}                                                   \tag{UL8}
\]
Every coefficient introduced on squaring the last two identities
is infinitesimal by UL2--UL3 and $S,u,v\in\mathrm I$.
Finally the formerly dangerous QI6 term has the exact replacement
\[
 \frac{wg}{c}\,vB_*C_*
       =wz\,C_*^2-\frac{wg}{c\sqrt\lambda}\,vC_*J_3,
 \qquad
 \frac{wg}{c\sqrt\lambda}
          =w\sqrt{\frac{z^2}{n}}=O(1).                          \tag{UL9}
\]
Both resulting coefficients are infinitesimal. This retains the
entire term for arbitrary independent source and testing rates.

\subsection{Both unrestricted state-normal corrections}

The new complete first $N$ coefficient, including the remainder
of TU17, differs from the formal generic SC coefficient by
\[
\begin{split}
\delta\mathcal L_N={}&
 8\lambda c(1-e)\mathsf F_N B_*
 +8\lambda q\widetilde d\,ubvg(Sa_tH-D_uv),\\
\mathsf F_N={}&
 \frac{P_bD_v-(1+qz)P_m}{P}-zD_uD_v
                 =V-zUD_v+qF_{Nq}.
\end{split}                                                     \tag{UL10}
\]
The last equality is a finite monic $q$ division; its coquotient
has bounded arrays because the numerator at $q=0$ is exactly
$V-zUD_v$. After division by $\sqrt{H_0}$, its $q$ part is
infinitesimal on $B_*$ by UL3. The other parts use
\[
\begin{split}
\lambda cV B_*&=\lambda qg\,vC_*
                            -\sqrt{\lambda q}\,v^2J_3,\\
\lambda czU B_*&=\lambda qbg\,v(uE_*+bC_*).
\end{split}                                                     \tag{UL11}
\]
Their coefficients on $\widehat N$ are infinitesimal since
$g/\sqrt{H_0}$ and $\sqrt{\lambda q/H_0}$ are bounded.
The $\widetilde d$ term in UL10 uses
$uH=R(uE_*+bC_*)$ and $uv=bC_*$; it has the same
infinitesimal coefficient bounds. Thus this is a lift of the
complete first normal, with no state component put to zero.

For the second normal it is convenient to compare the coefficient
after multiplication by $\sqrt\sigma$. Differentiating the full
generic SC3 identity, including the restored $D$ packet, gives
\[
 \sqrt\sigma\,\delta\mathcal L_y
       =8\lambda q\{\mathsf F_HH+\mathsf F_vv\},                 \tag{UL12}
\]
where
\[
\begin{split}
\mathsf F_H={}&D_v\{\rho_0\widetilde a
                    -(1-\rho_0)\widetilde c\}
 -D_v\widetilde aU\rho_0/D_u
 +\widetilde\alpha zS^2V\rho_0/D_u
 -\widetilde dS^2a_t,\\
\mathsf F_v={}&\widetilde dSD_u-\widetilde\alpha zS.
\end{split}
\]
These arrays are bounded through both source openings. Exactly,
\[
 \mathsf F_H|_{q=U=V=0}
       =\frac{(1-e)S}{1-S(1-h)},\qquad
 \mathsf F_v|_{q=U=V=0}
       =-\frac{(1-e)S}{1-S(1-h)}.                              \tag{UL13}
\]
As $R=1$ there, finite monic division gives
$R\mathsf F_H+\mathsf F_v=qF_q+UF_U+VF_V$ with bounded
arrays. Since $H=R(E_*+v)$, the whole UL12 coefficient is
therefore an infinitesimal combination of $E_*,F_*,A_*,C_*,D_*$:
use only
\[
 qv=cA_*-qE_*,\qquad Uv=buC_*,\qquad
 Vv=v(F_*+gD_*-vE_*).                                         \tag{UL14}
\]
The factor $1/\sqrt\sigma$ is bounded.
Changing the formal harmonic-normal normalization to the actual
one multiplies its previously bounded arrays by a pure-source
unit with residue one.

For completeness, the unchanged formal first $N$ coefficient is
bounded after the actual $\sqrt{H_0}$ normalization:
$h_2^0,h_b^0=O(q)$ and $h_b^0-h_2^0=O(qz)$, while
$g/\sqrt{H_0}$ and $c/\sqrt{H_0}$ are bounded. These are exactly
the two estimates used in the QI normal allocation.
All remaining hard coefficients from the original $C,A,D,\alpha$
lines have the same residues as before; the relative $C,A$
changes and the absolute $D,\alpha$ changes are infinitesimal.
Together with UL5 the complete hard block on
$\widehat N,N_y,J_3$ has a positive definite real residue.

\subsection{A finite contact lift for the additional row}

It remains to allocate the full scalar correction on roots that
vanish at contact. Merely bounding a coefficient on $A_*,B_*,C_*,D_*$
would not suffice. Keep the complete strict toric root list SC16,
the QI rows $E_*,F_*,R_F,R_E$, and add
\[
 V_3=\sqrt{\lambda\Gamma}(\sqrt{U_*}-\sqrt{B_*}).
\]
Since $U_*B_*=\Gamma^2$, the exact identities are
\[
\begin{split}
 J_3^2&=4V_3^2+\lambda(\sqrt{U_*}-\sqrt{B_*})^4,\\
 \Gamma J_3&=\frac{V_3^2}{\sqrt\lambda}
                      -2\sqrt\lambda\,\Gamma(B_*-\Gamma).
\end{split}                                                     \tag{UL15}
\]
Reserve half the $J_3^2$ coefficient in UL5 for the first
identity. It supplies strict real pivots on both $J_3,V_3$,
with the quartic square retained nonnegative. The normal
coupling in UL5 remains an infinitesimal hard Gram correction.

Here are the additional first-jet coefficient checks needed for
the second identity, so that its factor $\sqrt\lambda$ causes
no hidden loss. The only new scalar terms linear in $J_3$ in
UL8--UL9 have coefficients, up to bounded arrays, of the forms
\[
 \epsilon_B=-2\sqrt n\,v,\qquad
 \epsilon_T=\frac{wg}{c\sqrt\lambda}\,v .
                                                               \tag{UL16}
\]
Both have bounded finite testing difference arrays and are
infinitesimal. At the actual pure-source contact $v_c=cg$,
\[
 \sqrt\lambda\,\epsilon_B(u_c,v_c)=-2n\in\mathrm I,\qquad
 \sqrt\lambda\,\epsilon_T(u_c,v_c)=wz=e\in\mathrm I.             \tag{UL17}
\]
Multiplication by any of the bounded arrays produced in UL7
preserves these conclusions. No other scalar term introduces
$J_3$ linearly: the other five identities in UL8, and the
$\lambda q$ scalar corrections, use only the old QI rows.

Write a term $\epsilon J_3A_i$ as
$\epsilon J_3(A_i-\Gamma)+\epsilon\Gamma J_3$.
The first summand is already an infinitesimal quadratic on the
strict list. Divide $\epsilon-\epsilon(u_c,v_c)$ finitely by
$u-u_c,v-v_c$ and apply the exact SC18 conductors to
$\Gamma(u-u_c)$ and $\Gamma(v-v_c)$. This allocates the
nonconstant part as infinitesimal quadratic corrections involving
$J_3$. For the constant part use UL15. Its $V_3^2$
coefficient is infinitesimal, and its remaining coefficient on
$\Gamma(B_*-\Gamma)$ is infinitesimal by UL17.

There are also new bare-$\Gamma$ $E_*,F_*$ terms when the
first five identities in UL8 are expanded and their $C_*,D_*$
factors are written as $(A_i-\Gamma)+\Gamma$. Put
$L=\lambda q$ and $M_1=nS/z^2$, both infinitesimal.
The additional $E_*$ coefficients are bounded-array multiples
of $Lbu$ and $Lgv$; the additional $F_*$ coefficients are
bounded-array multiples of $Lg$ and $M_1g$.
At $u_c=bc,v_c=cg$ their required contact ratios are
\[
 \frac{Lbu_c}{c}=LS,\qquad
 \frac{Lgv_c}{c}=Lz,\qquad
 \frac{Lg}{g}=L,\qquad
 \frac{M_1g}{g}=M_1 .
                                                               \tag{UL18}
\]
All are infinitesimal, with bounded finite difference arrays
before these contact substitutions. The exact QI7 identities
therefore allocate these additional terms by the same finite
division as the old $E_*,F_*$ terms. This check is necessary:
the two new normal rows cannot be ignored when applying the
contact-jet argument.

After these explicit additional steps, the remainder has exactly
the QI form: a bounded infinitesimal quadratic on the listed
roots, the now-verified bare-$\Gamma$ $E_*,F_*$ couplings, and
a toric expression
$\Gamma\sum_i\mu_i(A_i-\Gamma)+\Gamma^2\sigma$.
The old $E_*,F_*$ contact coefficient ratios remain the bounded
ones in QI7--QI8: their unchanged $q=0$ arrays depend on bounded
$w$ and do not require $w/(c/g)$ bounded.
Every retained packet, including UL5 and UL15, vanishes at the
actual contact. The complete actual scalar consequently has zero
contact value and all three zero first jets. Apply the two
explicit finite divisions SC17--SC20. Their three-jet balance
makes the constant toric first-jet coefficients equal; SC19 and
the variance roots then allocate this common term. This puts
the entire scalar on a bounded Gram with strictly positive
real residue.

The full normal lifts UL10--UL14 have exact contact zero as well.
Their $J_3,E_*,F_*$ terms already vanish there. The sum of their
remaining four monomial coefficients is zero at contact, so one
finite SC18 division allocates them on the same strict list.
Every coefficient function in these steps depends only on the
testing $u,v$. No bound on the state ratio $Q_b/Z$ is used, and
the third contact jet is used only for its exact toric balance.

The strict bounded-Gram lemma now proves tensor SOS for the whole
original parent on the small-angle positive negative-core chart.
All its radicals are radicals of separated positive monomials;
every inverse is pure source or pure testing.

\subsection{Projective completion and scope}

The state, sign, zero-axis and real/imaginary-state arguments are
the same exact original-core arguments at the end of TU.
In particular $Z=0$ is an original positive restriction, and the
negative-core $Z=1$ chart forces $W+uY>0$, permitting the
separated radicals above. The theorem R2O covers every testing
angle chart outside $u,v$ both infinitesimal without any
condition on $\lambda$ or the native $a$. This proves UL1 for
all tests and states.

Combining this theorem with TU covers the entire original R2
sector with native $a$ unbounded: its nonnegative
$\lambda=n/(tB)$ is either bounded, including zero, or positive
unbounded. The complementary bounded-$a$ theorem remains a
separate input to the final R2 join. This local result alone is
not the universal tensor-SOS theorem or finite-lift existence.

\paragraph{Verification boundary.}
The focused exact checker verifies UL7, UL10, UL13, the full
generic first-normal derivatives before any graph restriction,
all six UL8 conductors, UL9, UL11, UL14--UL17 and the actual
contact of $J_3$. The native source comparisons are the separately
verified TU9 dependency. The source inequalities, bounded finite
arrays, finite contact allocation and projective join are the
analytical proof above, not consequences of assertion counts.
